\section{Fazit und Ausblick}

Dieses ``Independent Coursework II'' hat die Klassifikation anhand von erlernten, in den Netzwerkparametern vorhandenen Konzepten untersucht. Hier haben sich mehrere Attributierungs-Methoden als sinnvoll für das Auffinden erwiesen.
Bei der Untersuchung der Attributierungen fällt auf, dass besonders die mittleren und höheren Layer des untersuchten Netzwerks einige Konzepte enthalten, die ausschlaggebend für das Auffinden von Mitosen sind. Diese weisen hohe prozentuale Anteile über alle im jeweiligen Layer vorhandenen Konzepte auf. Dagegen scheinen in unteren Layern die Konzepte unwichtiger, jedoch werden über die Methoden hinweg die selben Encodings gemacht.
Auch die exemplarische Untersuchung der Heatmaps zeigt, dass untere Layer weniger wichtig für die Mitosen-Klassifikation - und mittlere dagegen entscheidend sind.

Für zukünftige Arbeit wäre es interessant zu untersuchen, wie sich Konzepte respektive eines Multi-Klassifikators verhalten. Da über ein Netzwerk verstreut die Filter auch für ähnliche Klassen sprechen können, müsste untersucht werden, wo und wie diese Konzepte welche Klassenprädiktionen bestärken.